\subsection{Workflow Graph Construction and Pre-Execution Validation}
\label{subsec:workflow-graph}
 
Before any workflow can be executed, the system must transform the raw JSON
payload sent by the frontend into a structured, in-memory graph, and then
verify that the graph is semantically correct. This process occurs in two
successive phases: \textit{graph construction} and \textit{pre-execution
validation}. Both phases complete before a single job is enqueued, ensuring
that the execution engine only receives well-formed, executable workflows.
 
\subsubsection{Data Model}
 
The type system that underpins the entire workflow engine is defined in
\texttt{workflow.type.ts}. Three core interfaces model the graph structure.
 
\textbf{\texttt{INode}} represents a single node in the workflow. Each node
carries all the information required to execute it: its unique identifier, the
node type name used to look it up in the registry, the project and execution
identifiers, an \texttt{inputSource} policy that decides where the node reads
its inputs from, the actual input and output data arrays, typed parameters,
credential references, and the list of its resolved children.
 
\begin{lstlisting}[caption={Core node interface (\texttt{workflow.type.ts})}]
export interface INode {
  id:           string;   // unique node id within the workflow
  type:         string;   // registry key, e.g. "cron", "poll"
  label:        string;   // human-readable display name
  projectId:    string;
  executionId:  string;   // assigned by the worker at runtime
  inputSource:  'self' | 'parent' | 'combined' | 'any';
  inputs:       string[]; // user-provided inputs for this node
  parentOutputs: string[];// outputs received from the upstream node
  parameters:   any[];
  credentials?: { typeId: string; setId: string }[];
  children:     INode[];  // resolved at deployment time
}
\end{lstlisting}
 
\textbf{\texttt{IConnection}} is deliberately minimal, capturing only the
directed edge between two node identifiers:
 
\begin{lstlisting}[caption={Connection and graph interfaces}]
export interface IConnection {
  source: string;
  target: string;
}
 
export interface IWorkflowGraph {
  nodes:       INode[];
  connections: IConnection[];
}
\end{lstlisting}
 
\textbf{\texttt{IWorkflowExecutionContext}} is the object passed from the API
controller to the workflow runner. It bundles the fully constructed graph with
the project identifier so that downstream functions do not need to accept
multiple unrelated arguments:
 
\begin{lstlisting}
export interface IWorkflowExecutionContext {
  workflowGraph: WorkflowGraph;
  projectId:     string;
}
\end{lstlisting}
 
\subsubsection{The \texttt{WorkflowGraph} Class}
 
The \texttt{WorkflowGraph} class wraps an \texttt{IWorkflowGraph} value
object with an \textit{adjacency list} for efficient traversal, and exposes
graph operations used by both the validation and execution layers.
 
\textbf{Internal Representation}
 
The class maintains two parallel data structures:
 
\begin{itemize}
  \item \texttt{graph} --- the plain \texttt{IWorkflowGraph} object, which
        stores the canonical lists of nodes and connections and is used
        wherever a serialisable snapshot of the graph is needed.
  \item \texttt{adjList} --- a \texttt{Map<string, string[]>} that maps each
        node identifier to the identifiers of its direct successors. This is
        kept in sync with \texttt{graph.connections} whenever
        \texttt{addConnection} is called, allowing $O(1)$ neighbour lookup
        during traversal instead of scanning the full connection list each time.
\end{itemize}
 
\begin{lstlisting}[caption={WorkflowGraph construction (\texttt{workflowGraph.class.ts})}]
export class WorkflowGraph {
  graph:           IWorkflowGraph;
  private adjList: Map<string, string[]>;
 
  constructor() {
    this.graph   = { nodes: [], connections: [] };
    this.adjList = new Map();
  }
 
  addNode(node: INode): void {
    if (this.hasNode(node))
      throw new AppError(`Node with ID ${node.id} already exists`, 400);
    this.graph.nodes.push(node);
  }
 
  addConnection(edge: IConnection): void {
    if (this.hasConnection(edge))
      throw new AppError(
        `Edge from ${edge.source} to ${edge.target} already exists`, 400);
 
    this.graph.connections.push(edge);
 
    if (!this.adjList.has(edge.source))
      this.adjList.set(edge.source, []);
    this.adjList.get(edge.source)!.push(edge.target);
  }
}
\end{lstlisting}
 
Both \texttt{addNode} and \texttt{addConnection} perform a duplicate check
before mutating state, throwing an \texttt{AppError} if the node or edge is
already present. This prevents malformed workflow payloads from producing
silent inconsistencies in the graph.
 
\textbf{Finding Trigger Nodes}
 
A trigger node is structurally defined as a node with no incoming edges ---
it sits at the root of a connected component. \texttt{findTriggerNodes}
identifies these nodes using an in-degree approach:
 
\begin{lstlisting}[caption={Trigger node detection by in-degree}]
findTriggerNodes(): INode[] {
  const allIDs      = this.getNodes().map(n => n.id);
  const hasIncoming = new Set<string>();
 
  // collect every node that appears as a target of some edge
  for (const node of this.getNodes())
    this.getNodeConnections(node.id).forEach(n => hasIncoming.add(n));
 
  // nodes absent from hasIncoming have no parents -> they are triggers
  const triggerIDs = allIDs.filter(id => !hasIncoming.has(id));
  return this.graph.nodes.filter(n => triggerIDs.includes(n.id));
}
\end{lstlisting}
 
The method builds the set of all nodes that appear as an edge target. Any
node \textit{not} in that set has an in-degree of zero and is therefore a
candidate trigger. The subsequent validation step then confirms, via the node
registry, that each such candidate is actually declared as a trigger node type.
 
\textbf{Decomposing into Connected Components}
 
A workflow canvas may contain multiple independent sub-workflows. Rather than
validating and executing the entire graph as one unit, the system first
decomposes it into \textit{connected components} so that each sub-workflow can
be validated and reported on individually. The decomposition uses BFS over the
\textit{undirected} version of the graph (edges are traversed in both
directions):
 
\begin{lstlisting}[caption={Connected component decomposition (BFS)}]
findConnectedComponents(): WorkflowGraph[] {
  const visited           = new Set<string>();
  const components: WorkflowGraph[] = [];
  const nodeMap           = new Map(this.graph.nodes.map(n => [n.id, n]));
  const reverseConnections = new Map<string, string[]>();
 
  // build a reverse adjacency map to traverse edges backwards
  for (const conn of this.graph.connections) {
    if (!reverseConnections.has(conn.target))
      reverseConnections.set(conn.target, []);
    reverseConnections.get(conn.target)!.push(conn.source);
  }
 
  for (const startNode of this.graph.nodes) {
    if (visited.has(startNode.id)) continue;
 
    const componentGraph = new WorkflowGraph();
    const queue          = [startNode];
    visited.add(startNode.id);
    componentGraph.addNode(startNode);
 
    while (queue.length > 0) {
      const current     = queue.shift()!;
      const forward     = this.adjList.get(current.id)  || [];
      const backward    = reverseConnections.get(current.id) || [];
      const neighbours  = [...new Set([...forward, ...backward])];
 
      for (const nId of neighbours) {
        if (!visited.has(nId)) {
          visited.add(nId);
          componentGraph.addNode(nodeMap.get(nId)!);
          queue.push(nodeMap.get(nId)!);
        }
      }
    }
 
    // copy only the edges that belong to this component
    for (const conn of this.graph.connections)
      if (componentGraph.hasNode({ id: conn.source } as INode) &&
          componentGraph.hasNode({ id: conn.target } as INode))
        componentGraph.addConnection(conn);
 
    components.push(componentGraph);
  }
 
  return components;
}
\end{lstlisting}
 
Treating the graph as undirected during decomposition is important. A directed
BFS from a trigger node would miss any nodes upstream of a cycle or any
disjoint nodes that happen to feed into the same component via reverse edges.
Using both the forward adjacency list and the reverse map guarantees that every
node is assigned to exactly one component regardless of edge direction.
 
\subsubsection{Graph Construction from the HTTP Request}
 
The \texttt{getWorkflowGraph} function in
\texttt{workflow-pre-execution.validations.ts} is the entry point for turning
a raw request payload into a typed \texttt{WorkflowGraph} instance. It is
called by the workflow controller immediately after the project record is
fetched from the database.
 
\textbf{Input Format}
 
The frontend serialises the workflow canvas as a JSON object containing a
\texttt{nodes} array and a \texttt{connections} array. Each node object
follows the React Flow schema, where the node's semantic type is stored inside
a nested \texttt{data.nodeName} field rather than at the top level:
 
\begin{lstlisting}[caption={Example workflow payload from the frontend}]
{
  "projectId": "abc-123",
  "nodes": [
    {
      "id": "cron-1761747201552",
      "data": { "nodeName": "cron", "label": "Scheduled Trigger" },
      "parameters": [
        { "title": "interval", "value": 60000 },
        { "title": "start",    "value": 1700000000000 }
      ],
      "credentials": []
    },
    {
      "id": "sheet-1761747208825",
      "data": { "nodeName": "googleSheet", "label": "Append to Sheet" },
      "parameters": [...],
      "credentials": [{ "typeId": "cb7c448b-...", "setId": "f3a1..." }]
    }
  ],
  "connections": [
    { "source": "cron-1761747201552",
      "target": "sheet-1761747208825", "type": "custom-edge" }
  ]
}
\end{lstlisting}
 
\textbf{Parsing and Validation}
 
The function first validates the shape of the incoming data, then maps each
raw node object to a typed \texttt{INode}, and finally populates the graph:
 
\begin{lstlisting}[caption={Graph construction from request body}]
export const getWorkflowGraph =
    (body: Record<string, unknown>): WorkflowGraph => {
 
  // --- shape checks -----------------------------------------------
  const nodesRaw = (body as any).nodes;
  if (!Array.isArray(nodesRaw) ||
      nodesRaw.some(n => typeof n.id !== 'string'))
    throw new AppError('Invalid or missing nodes', 400);
 
  const connectionsRaw = (body as any).connections;
  if (!Array.isArray(connectionsRaw))
    throw new AppError('Invalid or missing connections', 400);
 
  const isValidConnections = connectionsRaw.every(conn =>
    typeof conn.source === 'string' &&
    typeof conn.target === 'string' &&
    typeof conn.type   === 'string'
  );
  if (!isValidConnections)
    throw new AppError('Invalid connections', 400);
 
  // --- mapping -------------------------------------------------------
  const nodes: INode[] = nodesRaw.map(n => ({
    id:           n.id,
    type:         n.data?.nodeName,     // unwrap from React Flow schema
    label:        n.data?.label ?? '',
    projectId:    body.projectId as string,
    executionId:  '',                   // assigned later by the worker
    inputSource:  n.inputSource ?? 'parent',
    inputs:       n.inputs,
    parentOutputs: [],
    parameters:   n.parameters  ?? [],
    credentials:  n.credentials ?? [],
    children:     [],                   // resolved at deployment time
  }));
 
  const connections: IConnection[] = connectionsRaw.map(conn => ({
    source: conn.source,
    target: conn.target,
  }));
 
  // --- graph population ----------------------------------------------
  const graph = new WorkflowGraph();
  nodes.forEach(n => graph.addNode(n));
  connections.forEach(c =>
    graph.addConnection({ source: c.source, target: c.target }));
 
  return graph;
};
\end{lstlisting}
 
Two design details are worth noting. First, \texttt{executionId} is
intentionally left empty at this stage. The worker assigns it later when it
creates the \texttt{Execution} database record, and then propagates it through
all child jobs. Second, \texttt{children} is left as an empty array; the
workflow runner populates it recursively at deployment time via
\texttt{buildChildren}, so that the full execution tree travels with each job
without requiring further graph lookups.
 
\subsubsection{Pre-Execution Validation Pipeline}
 
Validation is orchestrated by \texttt{validateAllWorkflows}, which runs each
connected component through a sequence of checks. The function returns a
result object that separates valid sub-workflows from invalid ones, along with
the specific error messages for each failure:
 
\begin{lstlisting}[caption={Validation orchestrator}]
export const validateAllWorkflows =
    async (workflowGraph: WorkflowGraph) => {
 
  const validationResults = {
    validWorkflows:   [] as IWorkflowGraph[],
    invalidWorkflows: [] as { workflow: IWorkflowGraph;
                              errors: string[] }[],
  };
 
  const components = workflowGraph.findConnectedComponents();
 
  for (const component of components) {
    const errors: string[] = [];
 
    const unregisteredNodes = checkAllNodesValidation(component);
    const triggerResult     = validateTriggerNodes(component);
 
    if (!triggerResult.ok)
      errors.push(triggerResult.msg!);
 
    if (unregisteredNodes.length > 0) {
      errors.push(`The following nodes are not registered: ` +
        unregisteredNodes
          .map(n => `${n.name} (id: ${n.id})`).join(', '));
 
    } else if (triggerResult.ok && unregisteredNodes.length === 0) {
      // cycle check only runs when the graph structure is already valid
      const cycleResult = checkCycles(component);
      if (!cycleResult.ok)
        errors.push(cycleResult.msg!);
    }
 
    if (errors.length > 0)
      validationResults.invalidWorkflows.push(
        { workflow: component.graph, errors });
    else
      validationResults.validWorkflows.push(component.graph);
  }
 
  return validationResults;
};
\end{lstlisting}
 
The three validators are applied in a deliberate order, described in the
following subsubsections.
 
\textbf{Trigger Node Validation}
 
\texttt{validateTriggerNodes} enforces three structural rules:
 
\begin{enumerate}
  \item The component must contain at least one node.
  \item Every node with no incoming edges (as identified by
        \texttt{findTriggerNodes}) must be declared as a trigger node in the
        registry (\texttt{triggerNode: true}).
  \item No node with incoming edges may be a trigger node --- trigger nodes
        are only valid at the root of a component.
\end{enumerate}
 
\begin{lstlisting}[caption={Trigger node structural validation}]
const validateTriggerNodes =
    (graph: WorkflowGraph): { ok: boolean; msg?: string } => {
 
  const triggerNodes = graph.findTriggerNodes();
 
  if (graph.getNodes().length === 0)
    return { ok: false, msg: 'No nodes found to execute.' };
 
  if (triggerNodes.length === 0)
    return { ok: false, msg: 'No start node found or cycle exists.' };
 
  // every root must be a registered trigger
  for (const node of triggerNodes) {
    const desc = registry.getDescription(node.type);  // throws if unregistered
    if (!desc.triggerNode)
      return { ok: false,
        msg: `Node '${node.type}' (id: '${node.id}') is at the ` +
             `start of the workflow but is not a trigger node.` };
  }
 
  // no non-root node may be a trigger
  for (const node of graph.getNodes()) {
    if (triggerNodes.includes(node)) continue;
    const desc = registry.getDescription(node.type);
    if (desc.triggerNode)
      return { ok: false,
        msg: `Trigger node '${node.type}' (id: '${node.id}') ` +
             `cannot be used in the middle of a workflow.` };
  }
 
  return { ok: true };
};
\end{lstlisting}
 
\textbf{Node Registry Validation}
 
\texttt{checkAllNodesValidation} ensures that every non-trigger node in the
component is registered in the \texttt{NodeRegistry}. Trigger nodes are
excluded from this check because their registration is already verified by
\texttt{validateTriggerNodes}.
 
\begin{lstlisting}[caption={Registry membership check for non-trigger nodes}]
const checkAllNodesValidation =
    (workflowGraph: WorkflowGraph): { name: string; id: string }[] => {
 
  const triggerNodeIds = new Set(
    workflowGraph.findTriggerNodes().map(n => n.id));
  const unregistered: { name: string; id: string }[] = [];
 
  for (const node of workflowGraph.graph.nodes) {
    if (triggerNodeIds.has(node.id)) continue;  // skip triggers
    if (!registry.has(node.type))
      unregistered.push({ name: node.type, id: node.id });
  }
 
  return unregistered;
};
\end{lstlisting}
 
If any unregistered nodes are found, their names and identifiers are collected
and returned as a list, which the orchestrator formats into a single error
message. This allows the user to see all unregistered nodes in one response
rather than having to fix them one by one.
 
\textbf{Cycle Detection}
 
The cycle check is the last to run and only executes if both preceding checks
pass. This ordering is intentional: running a DFS-based cycle detector on a
graph that contains unregistered nodes or a missing trigger would produce
unreliable or misleading results, since the graph structure cannot be trusted.
 
Cycle detection uses a standard three-colour DFS algorithm, encoded by the
\texttt{VisitState} enum:
 
\begin{lstlisting}[caption={\texttt{VisitState} enum (\texttt{workflow.enum.ts})}]
export enum VisitState {
  UNVISITED = 0,  // not yet reached
  VISITING  = 1,  // currently on the DFS call stack
  VISITED   = 2,  // fully processed, no cycle through this node
}
\end{lstlisting}
 
\begin{lstlisting}[caption={DFS cycle detection}]
const checkCycles =
    (graph: WorkflowGraph): { ok: boolean; msg?: string } => {
 
  const visited = new Map<string, VisitState>(
    graph.getNodes().map(n => [n.id, VisitState.UNVISITED]));
 
  const dfs = (u: string): boolean => {
    visited.set(u, VisitState.VISITING);
 
    for (const v of graph.getNodeConnections(u)) {
      if (visited.get(v) === VisitState.VISITING) return true;  // back edge
      if (visited.get(v) === VisitState.UNVISITED && dfs(v))   return true;
    }
 
    visited.set(u, VisitState.VISITED);
    return false;
  };
 
  for (const node of graph.getNodes())
    if (visited.get(node.id) === VisitState.UNVISITED)
      if (dfs(node.id))
        return { ok: false, msg: 'Workflow has cycles' };
 
  return { ok: true };
};
\end{lstlisting}
 
A node in the \texttt{VISITING} state is one that is currently on the
recursive call stack. Encountering such a node while traversing its own
descendants means a back edge has been found, which is the defining
characteristic of a cycle in a directed graph. Once all descendants of a node
have been processed without finding a back edge, the node is marked
\texttt{VISITED} and will not be re-explored. The outer loop ensures that all
disconnected components --- if \texttt{findConnectedComponents} was bypassed
--- are also checked.
 
\subsubsection{Validation Ordering and Error Accumulation}
 
The order in which the three validators run is not arbitrary. Table
\ref{tab:validation-order} summarises the rationale for their sequence.
 
\begin{table}[h!]
\centering
\begin{tabular}{clp{7.5cm}}
  \toprule
  \textbf{Order} & \textbf{Validator} & \textbf{Rationale} \\
  \midrule
  1 & \texttt{validateTriggerNodes}   &
      Structural check. A component without a valid trigger root cannot be
      executed at all, and a trigger node in the middle of a graph is a
      logical error that makes the remaining checks meaningless. \\[6pt]
  2 & \texttt{checkAllNodesValidation} &
      Registry check. An unregistered node would cause the worker to throw at
      runtime. Catching this statically produces a clearer error and avoids
      wasting a job queue slot. \\[6pt]
  3 & \texttt{checkCycles} &
      Graph-theoretic check. Only runs when the graph is structurally sound.
      A cyclic graph would cause \texttt{buildChildren} to recurse infinitely
      and flood the job queue, so cycles must be rejected before deployment. \\
  \bottomrule
\end{tabular}
\caption{Pre-execution validation order and rationale}
\label{tab:validation-order}
\end{table}
 
Errors from the trigger and registry checks are accumulated independently and
both reported if both fail. The cycle check is \textit{gated} behind the
registry check: it only executes if all non-trigger nodes are registered, since
a DFS over a graph that references unknown node types cannot be trusted to
reflect the true structure of the intended workflow.
 
Once the full validation result is returned to the workflow controller, any
component listed in \texttt{invalidWorkflows} causes the entire deployment
request to be rejected with a structured error payload identifying the
problematic nodes and the reasons for failure.
 
